\providecommand{\kBKM}{\kappa_{\mathrm{BKM}}}
\providecommand{\kch}{\kappa_{\mathrm{ch}}}
\providecommand{\kcat}{\kappa_{\mathrm{cat}}}
\providecommand{\kfib}{\kappa_{\mathrm{fiber}}}
\providecommand{\Hilb}{\mathrm{Hilb}}
\providecommand{\DT}{\mathrm{DT}}
\providecommand{\cP}{\mathcal{P}}
\providecommand{\cC}{\mathcal{C}}
\providecommand{\HH}{\mathbb{H}}
\providecommand{\Z}{\mathbb{Z}}
\providecommand{\Pone}{\mathbb{P}^{1}}
\providecommand{\CoHA}{\mathrm{CoHA}}
\providecommand{\SpCh}{\mathrm{Sp}^{\mathrm{ch}}}
\providecommand{\PhiFA}{\Phi^{\mathrm{FA}}}

\section*{Automorphic scalar and Hall gluing data}

The automorphic calculation fixes an arithmetic skeleton: the
hyperbolic lattice $\Lambda^{2,1}$, its reflection polyhedra, the
Borcherds lift $\phi_{0,1}\mapsto\Delta_5$, and the reduced
Donaldson--Thomas scalar attached to
\[
 X=(S\times E)/(\Z/N\Z).
\]
A geometric gluing theorem requires further data: an open or stacky
atlas $\{U_i\}$, \v{C}ech cocycles on overlaps, local quasi-NCCR
algebras or their replacements on strata, transition equivalences, and
a Behrend-compatible descent map from the local Hall models to the
reduced $K3\times E$ Donaldson--Thomas integral. The automorphic
calculation supplies lattice and scalar constraints. The gluing theorem
starts from those constraints and constructs the missing geometry.

\paragraph{Proposition 1 (reflection chambers live in the lattice cone).}
The polyhedra
\[
\cP,\qquad \cP_I,\qquad \cP_{II}
\]
cut the cone
$\cC(\Lambda^{2,1})_+\subset\Lambda^{2,1}\otimes\mathbb R$ by
reflection hyperplanes. Their stabilisers are
\[
A(\cP)=\langle 1\rangle,\qquad
A(\cP_I)=\langle s_{f_3}\rangle\cong S_2,\qquad
A(\cP_{II})=\langle s_{f_2-f_3},s_{f_{-2}-f_2}\rangle\cong S_3,
\]
and
\[
O(\Lambda^{2,1})_+\cong W^{(2)}(\Lambda^{2,1}_{II})\rtimes S_3
\cong \mathrm{PGL}_2(\Z).
\]
This is a Weyl-chamber tiling of a hyperbolic cone. It is separate
from a \v{C}ech cover of $K3\times E$, a toric atlas of
$(S\times E)/(\Z/N\Z)$, and an NCCR gluing diagram.

\emph{Proof.}
The defining data are primitive square-two roots
\[
\cP'=\bigcap_{\delta\in\cP'_{\mathrm{prim}}}\mathcal H^+_\delta
\]
inside $\Lambda^{2,1}$. Reflections $s_\delta$ act on the lattice cone;
their semidirect product records the automorphy of the Borcherds
denominator under the Weyl group. A \v{C}ech cover of a scheme or
stack consists of morphisms $U_i\to X$, overlap schemes
$U_i\times_X U_j$, and cocycles on triple overlaps. The reflection
calculation names roots, walls, chambers, and Weyl symmetries. It
stops before schemes $U_i$, overlap morphisms, algebras on charts, and
transition functors. Local geometric descent for $X$ requires the atlas
and transition data named above.

\paragraph{Proposition 2 (the DT block is scalar trace data).}
The quotient $X=(S\times E)/(\Z/N\Z)$ comes with an elliptic K3 datum
$\pi:S\to\Pone$, torsion sections $s_1,\ldots,s_N$, and classes
\[
\beta_h=\frac{s_1+\cdots+s_N+hF}{N}.
\]
The reduced Behrend-weighted series is
\[
Z^X(q,t,p)
=\sum_{h\geq0}\sum_{d\geq0}\sum_{n\in\Z}
\DT^{X,\mathrm{red}}_{n,(\beta_h,d)}\,q^{d-1}
t^{\frac12\langle\beta_h,\beta_h\rangle}(-p)^n,
\]
with
\[
\DT^{X,\mathrm{red}}_{n,(\beta_h,d)}
=\sum_{k\in\Z}k\,e\!\left(\nu^{-1}(k)\right)
\quad\text{on}\quad
\Hilb^n(X,(\beta_h,d))/E .
\]
In the Oberdieck--Pandharipande monic normalisation,
\[
D_5=64^{-1}\Delta_5,\qquad
\Phi_{10}^{\mathrm{OP}}=D_5^2=4096^{-1}\Delta_5^2,
\]
and at $N=1$
\[
Z^{K3\times E}_{\DT,\mathrm{red}}
=-\bigl(\Phi_{10}^{\mathrm{OP}}\bigr)^{-1}
=-D_5^{-2}
=-4096\,\Delta_5^{-2}.
\]

\emph{Proof.}
The Behrend function is integrated over the Hilbert scheme modulo the
elliptic translation action. The quotient by $E$ removes the trivial
translation direction in the invariant integral and produces the
reduced scalar trace. The equality
$4096=64^2$ is the square of the monic normalisation
$D_5=64^{-1}\Delta_5$. A descent system names local
algebras $\mathcal A_i$, equivalences
$\mathcal A_i|_{U_{ij}}\simeq\mathcal A_j|_{U_{ij}}$, cocycle
compatibilities on triple overlaps, and a character map from the glued
category to reduced DT integration. The scalar trace gives the value
that such a construction must realise.

\paragraph{Proposition 3 (Borcherds constants are input coefficients).}
For the CHL ladder $N\in\{1,2,3,4,6\}$ the Borcherds-input constants and
BKM modular characteristics are
\[
(c_1(0),c_2(0),c_3(0),c_4(0),c_6(0))=(10,8,6,4,2),
\]
\[
(\kBKM(\Phi_1),\kBKM(\Phi_2),\kBKM(\Phi_3),
\kBKM(\Phi_4),\kBKM(\Phi_6))=(5,4,3,2,1),
\]
with
\[
\kBKM(\Phi_N)=\frac{c_N(0)}2 .
\]
At $N=1$,
\[
\kBKM(\Delta_5)=\frac{c_1(0)}2=\frac{10}{2}=5.
\]

\emph{Proof.}
Borcherds' singular-theta weight formula (Borcherds 1998,
Theorem~13.3) identifies the weight of the denominator product with
half the constant coefficient of the input Jacobi form. The K3 elliptic
genus has
$\phi_{0,1}|_{q^0}=\zeta^{-1}+10+\zeta$, giving the $N=1$ value.
The Gritsenko--Nikulin CHL forms (Gritsenko--Nikulin 1995, Part~II,
Theorem~2.1; Gritsenko 1999, Theorem~1.2) have weights
$(5,4,3,2,1)$ on $N=(1,2,3,4,6)$. The constants are twice the weights.
These are BKM weights. Hodge supertraces and fibre Euler
characteristics are separate invariants.

\paragraph{Proposition 4 (the $K3\times E$ invariants are separated).}
For $Y=K3\times E$, the compact total-space values are
\[
\kcat(Y)=\chi(\mathcal O_{K3})\chi(\mathcal O_E)=2\cdot0=0,
\]
\[
\kch(Y)=\sum_q(-1)^q h^{0,q}(Y)=1-1+1-1=0.
\]
The three additional branch values are
\[
\kch^{\mathrm{Heis}}(Y)=3,\qquad
\kBKM(\Delta_5)=5,\qquad
\kfib(K3)=\mathrm{rk}\,\widetilde H(K3,\Z)=24.
\]
Thus the quoted spectrum
\[
\{0,3,5,24\}
\]
records the compact zero together with the Heisenberg--Mukai
specialisation, the Borcherds weight, and the Mukai-lattice rank.

\emph{Proof.}
K\"unneth multiplicativity gives $\chi(\mathcal O_{K3\times E})=0$.
The compact chiral Hodge supertrace on the CY$_3$ total space also
vanishes. The value $3$ is read from the Stage-$2$
Heisenberg--Mukai specialisation
$\SpCh_{\Sigma_2,E}\circ\PhiFA_3$ with $\Sigma_2=K3$ and target curve
$E$. The value $5$ is the Borcherds weight of $\Delta_5$. The value
$24$ is the Mukai rank $1+22+1$ (Mukai 1987). The fibre witness
$\chi(\mathcal O_{K3})=2$ has its own role and is distinct from the
total-space value. The additive formula
\[
\kBKM=\kch+\chi(\mathcal O_{\mathrm{fiber}})
\]
fails already at $N=1$, where the left side is $5$ and the compact
right side is $0+0$. The CHL family is governed by
$\kBKM(\Phi_N)=c_N(0)/2$.

\paragraph{Proposition 5 (construction routes and the single \(\Phi\)
output).}
The six routes to $G(K3\times E)$ are six construction procedures:
Schiffmann--Vasserot CoHA, Maulik--Okounkov stable envelopes,
Gritsenko--Nikulin Borcherds lift, Costello--Paquette holographic
input, Bridgeland--Smith autoequivalence geometry, and
Fourier--Jacobi/Gopakumar--Vafa genus escalation. The functorial input
is the single Stage-$1$ object
\[
\PhiFA_3(\mathrm{Perf}(K3\times E)).
\]
At $d\geq3$ the native specialisation target of the CY-to-chiral
construction is $E_1$-chiral. The $E_2$ structure belongs to the
derived centre, a Hall--Drinfeld double, or another constructed
comparison object.

\emph{Proof.}
The Stage-$1$ object has many Stage-$2$ specialisations because the
choice of $(\Sigma_2,C)$ varies. Comparison with a common completed
target requires the Hall package. This is a multiplicity of construction
routes from one CY$_3$ input. The Dunn--Lurie dimension drop fixes the
native $d=3$ specialised output as $E_1$-chiral; the braided $E_2$
object is
$\mathcal Z(\Rep^{E_1}(A))$ or the corresponding completed
Hall--Drinfeld double after the centre, pairing, and completion have
been constructed.

\paragraph{Proposition 6 (positive Hall data precede doubles and
centres).}
The local toric normalisation is
\[
\CoHA(\mathbb C^3)=Y^+(\widehat{\mathfrak{gl}}_1),
\]
the positive half of the affine Yangian. The full comparison passes
through
\[
Y^+(\widehat{\mathfrak{gl}}_1)
\hookrightarrow
D\bigl(Y^+(\widehat{\mathfrak{gl}}_1)\bigr)
=Y^+\otimes Y^0\otimes Y^-,
\]
and then through the relevant centre, completion, or vertex evaluation.
For compact $K3\times E$, the gluing theorem must first construct the
positive Hall object, its product, orientation data, and candidate
pairing before taking a Drinfeld double or identifying a BKM bracket.

\emph{Proof.}
The CoHA product is a convolution product over effective curve classes.
It naturally sees the positive cone. The negative half, Cartan half,
Hopf pairing, completion topology, and centre are additional algebraic
structures. A scalar equality such as
$Z^{K3\times E}_{\DT,\mathrm{red}}=-4096\,\Delta_5^{-2}$ supplies a
trace constraint on a completed object. The Hall product, negative half,
centre, and bracket-recognition map are additional algebraic data.

\paragraph{Proposition 7 (order-six local ADE type and Niemeier anchor).}
For an order-six symplectic K3 automorphism, the cyclic quotient surface
has local Du Val type
\[
2A_5+2A_2+2A_1 .
\]
The selected CHL/umbral Niemeier anchor is the rank-$24$ root system
\[
6D_4 .
\]
The root system $4A_5+D_4$ belongs to a different Coxeter-six Niemeier
row.

\emph{Proof.}
The local ADE type is computed from stabiliser orbits of the cyclic
action on the K3 surface: two full $\Z/6$ fixed points give two
$A_5$ singularities, the remaining $\Z/3$ fixed orbits give two
$A_2$ singularities, and the remaining $\Z/2$ fixed orbits give two
$A_1$ singularities. The Niemeier label records a global even
unimodular rank-$24$ lattice selected by the umbral class
(Mukai 1988; Cheng--Duncan--Harvey 2014). Local surface singularities
and global Niemeier anchors live in different categories.

\paragraph{Gluing data required.}
A theorem upgrading the automorphic scalar to positive geometry must
construct the following data.
\begin{enumerate}
\item An open or stacky atlas of $X=(S\times E)/(\Z/N\Z)$, with
overlaps and triple overlaps.
\item Local toric or analytic models on each stratum, including the
order-$N$ quotient charts.
\item A Serre-equivariant quasi-NCCR or replacement algebra on the
admissible strata, with transition equivalences on overlaps.
\item A Behrend-compatible character map from the glued category to the
reduced DT integration.
\item Compatibility with the OP scalar
\[
Z^{K3\times E}_{\DT,\mathrm{red}}
=-\bigl(\Phi_{10}^{\mathrm{OP}}\bigr)^{-1}
=-D_5^{-2}
=-4096\,\Delta_5^{-2}
\]
at $N=1$, and with the CHL constants
\[
\kBKM(\Phi_N)=c_N(0)/2=(5,4,3,2,1),
\qquad N\in\{1,2,3,4,6\}.
\]
\item A positive Hall algebra on the effective cone, followed by the
pairing, completion, Drinfeld double or centre, and bracket-recognition
map to the intended Borcherds--Kac--Moody object.
\end{enumerate}

The automorphic calculation supplies lattice chambers, Borcherds
weights, and the OP scalar trace. A geometric gluing theorem begins
with the atlas, local models, quasi-NCCR transitions, reduced character
map, positive Hall algebra, pairing, completion, and BKM recognition
map.
